\section{Spinor Subalgebras and Kinematic Decoupling}

Before demonstrating confinement, we must identify the regime where standard associative kinematics (e.g., the Dirac equation) can operate. The split-quaternions $\mathbb{H}_s \cong M_2(\mathbb{R})$ form a 4-dimensional associative subalgebra within $\mathbb{O}_s$. 

By restricting the off-diagonal vectors to a single collinear axis (e.g., the first coordinate), the non-associative cross products completely vanish. We formalized this in Lean 4 by defining the property \texttt{isSplitQuaternion(X)} which enforces $\mathbf{u} = [u_0, 0, 0]$ and $\mathbf{v} = [v_0, 0, 0]$. Under this restriction, the algebra is proven to be strictly associative, recovering the $SL(2, \mathbb{R})$ spacetime spinor algebra. This guarantees that spacetime kinematics are strictly decoupled from the internal gauge sector.

\section{Gauge Automorphisms and $SU(3)$ Color}

The internal symmetry acting on the off-diagonal vectors $\mathbf{u}$ and $\mathbf{v}$ corresponds to the strong-force color charges. Any linear transformation $U: \mathbb{Z}^3 \to \mathbb{Z}^3$ that preserves the 3D dot product and cross product acts as a strict automorphism on the entire Zorn algebra. 

We identify this automorphism group as the $SU(3)$ color gauge group. Geometrically, $SU(3)$ emerges naturally as the stabilizer of the Peirce diagonal projections ($P_+$ and $P_-$) inside the exceptional group $G_2(2)$, leaving the associative spacetime observables invariant while rotating the internal color states.

\section{Derivation of the Hadronic Spectrum}

We now present the central result: the algebraic proof of hadronic confinement. Colored quarks and anti-quarks are defined as upper and lower nilpotent matrices, respectively:
\begin{equation}
Q_i = \begin{pmatrix} 0 & \mathbf{u}_i \\ \mathbf{0} & 0 \end{pmatrix}, \quad \bar{Q}_i = \begin{pmatrix} 0 & \mathbf{0} \\ \mathbf{v}_i & 0 \end{pmatrix}
\end{equation}
Because their Zorn norms vanish ($\text{zornNorm}(Q_i) = 0$), these individual states are strictly confined to the ultrahyperbolic light-cone. Furthermore, any arbitrary combination of these nilpotents yields off-diagonal vector components, meaning they are fundamentally non-associative and gauge-dependent (unobservable). 

However, specific topological combinations collapse exactly onto the diagonal, stripping them of all off-diagonal non-associativity.

\subsection{The Baryon Associator}
Using the Lean 4 proof assistant, we verified the exact Zorn algebra evaluation for the associator of three quarks $[Q_1, Q_2, Q_3] = (Q_1 \cdot Q_2) \cdot Q_3 - Q_1 \cdot (Q_2 \cdot Q_3)$. The result precisely mirrors the symbolic SymPy calculation:

\begin{equation}
B = [Q_1, Q_2, Q_3] = -T \begin{pmatrix} 1 & \mathbf{0} \\ \mathbf{0} & -1 \end{pmatrix} = -T \ell
\end{equation}
where $T = (\mathbf{u}_1 \times \mathbf{u}_2) \cdot \mathbf{u}_3$ is the scalar triple product (the $SU(3)$ color determinant). The baryon $B$ is completely stripped of off-diagonal vector components, making it associative and physically observable. Crucially, the baryon does not map to the scalar identity $I$, but rather to the paracomplex split-unit $\ell$. This identifies the baryon as an axial/paracomplex scalar, natively encoding chiral parity.

\subsection{The Anti-Baryon Associator}
By geometric symmetry, the associator of three anti-quarks yields a mirror paracomplex result:
\begin{equation}
\bar{B} = [\bar{Q}_1, \bar{Q}_2, \bar{Q}_3] = -\bar{T} \ell
\end{equation}
where $\bar{T} = (\mathbf{v}_1 \times \mathbf{v}_2) \cdot \mathbf{v}_3$. Baryons and anti-baryons share identical paracomplex stability, locked to the same grading axis $\ell$.

\subsection{The Meson Spectrum}
Combining a quark and an anti-quark allows for two distinct associative invariants depending on the choice of Jordan product:

\textbf{1. The Symmetric Meson (Real Scalar):}
The anticommutator $\{Q, \bar{Q}\} = Q \cdot \bar{Q} + \bar{Q} \cdot Q$ yields:
\begin{equation}
M_I = \{Q, \bar{Q}\} = (\mathbf{u} \cdot \mathbf{v}) \begin{pmatrix} 1 & \mathbf{0} \\ \mathbf{0} & 1 \end{pmatrix} = (\mathbf{u} \cdot \mathbf{v}) I
\end{equation}
This projects the meson onto the identity $I$, representing a standard associative particle (e.g., a pion) propagating through the real sector of the manifold.

\textbf{2. The Antisymmetric Meson (Paracomplex Scalar):}
The commutator $[Q, \bar{Q}] = Q \cdot \bar{Q} - \bar{Q} \cdot Q$ yields:
\begin{equation}
M_\ell = [Q, \bar{Q}] = (\mathbf{u} \cdot \mathbf{v}) \begin{pmatrix} 1 & \mathbf{0} \\ \mathbf{0} & -1 \end{pmatrix} = (\mathbf{u} \cdot \mathbf{v}) \ell
\end{equation}
This projects the meson onto the split-unit $\ell$, representing an axial-vector or paracomplex meson geometrically coupled to the thermal waves of the vacuum.

\end{document}
